TorchDCM is organized around a lightweight package layer and a separate validation layer.
The package layer contains installable code under \texttt{torchdcm/}, including data containers, specification objects, model classes, inference utilities, and result objects.
The validation layer contains public data processing, reference-estimator wrappers, benchmark scripts, and generated reports.
This organization keeps the user-facing package compact while preserving the ability to run heavyweight comparisons against Biogeme, Apollo, and R packages.

The central user-facing abstraction is a choice dataset.
For standard choice models, a \texttt{ChoiceDataset} stores alternatives, chosen alternatives, variables, availability masks, optional weights, and optional individual identifiers for panel likelihoods.
The package supports wide-to-long conversion for common transportation data layouts and ragged long-format choice sets when different observations have different feasible alternatives.
For ordered response models, an \texttt{OrderedChoiceDataset} stores ordinal outcomes, covariates, category labels, and optional observation weights.

Model specification is expressed through utility objects and named parameters.
Users define \texttt{Beta} parameters and combine them with observed variables in a \texttt{UtilitySpec}.
This syntax keeps the model readable while preserving a tensor-executable representation of each utility expression.
Advanced model objects extend the same specification layer: nested logit adds nests and dissimilarity parameters; cross-nested logit adds allocation weights; mixed logit adds random coefficients and simulation draws; ordered models add thresholds; and latent class models add class-specific utilities and membership models.

Result objects are designed to expose econometric outputs rather than only fitted tensors.
A fitted model returns parameter estimates, log likelihoods, covariance estimates, standard errors, t-statistics, probability predictions, and post-estimation summaries when available.
For MNL and related models, the package also reports willingness-to-pay and elasticity calculations used in the benchmark suite.
This reporting design is important because the package is evaluated not only by final likelihood values but also by whether downstream econometric quantities match reference implementations.
